% ============================================================
% preamble: packages & commands, shared by every version.
% (venue templates load many of these themselves; venue wrappers
% may skip this file & use only the "commands" part at the bottom.)
% ============================================================

\usepackage[T1]{fontenc}
\usepackage[utf8]{inputenc}
\usepackage{lmodern}
\usepackage{microtype}
\usepackage[letterpaper, margin=1in]{geometry}

\usepackage{amsmath, amssymb}
\usepackage{graphicx}
\usepackage{booktabs}
\usepackage{tabularx}
\usepackage{enumitem}
\usepackage[font=small, labelfont=bf]{caption}
\usepackage{subcaption}
\usepackage[dvipsnames]{xcolor}
\usepackage[numbers, sort&compress]{natbib}
\usepackage[hidelinks]{hyperref}

\graphicspath{{figures/}}

% ------------------------------------------------------------
% commands (used in the sections; keep these in every version)
% ------------------------------------------------------------

% a note to the author: shown while drafting, hidden for submission.
\newcommand{\note}[1]{\ifshownotes{\color{Magenta}\small[\textbf{note:} #1]}\fi}

% a placeholder for the author's own anecdote.
\newcommand{\anecdote}[1]{\ifshownotes\par\noindent{\color{Magenta}\small\fbox{\parbox{0.96\linewidth}{\textbf{anecdote (to write):} #1}}}\par\fi}

% text that names the author or links to their work; hidden when anonymous.
\newcommand{\identifying}[2]{\ifanonymous #2\else #1\fi}

% code & names from the system (kept in lowercase, as written).
\newcommand{\code}[1]{\texttt{#1}}

% the title of an interpretation (kept in lowercase, as the author writes them).
\newcommand{\work}[1]{\emph{#1}}

% the three parts of an interpretation, as written in its header.
\newenvironment{interpretation}[1]
  {\par\medskip\noindent\textbf{interpretation #1}\par\smallskip
   \begin{description}[leftmargin=1.2em, labelindent=0em, font=\normalfont\itshape, itemsep=0.25em]}
  {\end{description}\medskip}
